% =========================================================
% Algebra of Functions — Notes
% Chapter: functions
% Section: algebra-of-functions
% Volume III · Analysis
% =========================================================

\section{Algebra of Functions}
\label{sec:algebra-of-functions}

\begin{tcolorbox}[title={Toolkit --- Algebra of Functions}]
\begin{itemize}
  \item All results are set-theoretic: no limits, no topology.
  \item $\operatorname{Injective}(f)\wedge\operatorname{Injective}(g)
    \Rightarrow\operatorname{Injective}(g\circ f)$.
  \item $\operatorname{Bijective}(f)\Rightarrow\operatorname{Bijective}(f^{-1})$.
  \item Preimage: $f^{-1}(A\cup B)=f^{-1}(A)\cup f^{-1}(B)$,
    $f^{-1}(A\cap B)=f^{-1}(A)\cap f^{-1}(B)$ --- exact.
  \item Image: $f(A\cup B)=f(A)\cup f(B)$ but
    $f(A\cap B)\subseteq f(A)\cap f(B)$ --- only $\subseteq$ in general.
\end{itemize}
\end{tcolorbox}

\begin{remark*}[Intervals on $\mathbb{R}$]
Every open interval $(a,b)$ consists entirely of cluster points;
$[a,b]$ is closed and bounded. Natural domains for functions are
typically intervals or unions of intervals.
\end{remark*}

% ---------------------------------------------------------
% Proposition — Composition of Injections
% ---------------------------------------------------------

\begin{proposition}[Composition of Injections Is Injective]
\label{prop:composition-injective}
Let $f:A\to B$ and let $g:B\to C$. If $f$ is injective and $g$ is
injective, then $g\circ f:A\to C$ is injective.

\smallskip\noindent
\hyperref[prf:composition-injective]{\textit{Go to proof.}}
\end{proposition}

\begin{remark*}[Standard quantified statement]
\[
\begin{aligned}
&\Bigl(\forall a_1,a_2\in A\;(f(a_1)=f(a_2)\Rightarrow a_1=a_2)\Bigr)
\wedge
\Bigl(\forall b_1,b_2\in B\;(g(b_1)=g(b_2)\Rightarrow b_1=b_2)\Bigr) \\
&\quad\Rightarrow
\forall a_1,a_2\in A\;(g(f(a_1))=g(f(a_2))\Rightarrow a_1=a_2).
\end{aligned}
\]
\end{remark*}

\begin{remark*}[Definition predicate reading]
\[
  \operatorname{Injective}(f)\wedge\operatorname{Injective}(g)
  \;\Longrightarrow\;
  \operatorname{Injective}(g\circ f).
\]
\end{remark*}

\begin{remark*}[Negated quantified statement]
\[
  \exists a_1,a_2\in A\;\bigl(a_1\neq a_2\;\wedge\;g(f(a_1))=g(f(a_2))\bigr).
\]
\end{remark*}

\begin{remark*}[Negation predicate reading]
\[
  \neg\,\operatorname{Injective}(g\circ f).
\]
\end{remark*}

\begin{remark*}[Failure modes]
The composite fails to be injective exactly when two distinct inputs of
$A$ are sent to the same output in $C$ after applying both functions.
\end{remark*}

\begin{remark*}[Failure mode decomposition]
\[
  \exists a_1,a_2\in A\;\bigl(a_1\neq a_2\;\wedge\;g(f(a_1))=g(f(a_2))\bigr).
\]
\end{remark*}

\begin{remark*}[Interpretation]
Injectivity survives composition: equality after $g\circ f$ can be
pulled back first through $g$ and then through $f$.
\end{remark*}

\NoLocalDependencies
% ---------------------------------------------------------
% Proposition — Composition of Surjections
% ---------------------------------------------------------

\begin{proposition}[Composition of Surjections Is Surjective]
\label{prop:composition-surjective}
$\operatorname{Surjective}(f)\wedge\operatorname{Surjective}(g)
\Rightarrow\operatorname{Surjective}(g\circ f)$.

\smallskip\noindent
\hyperref[prf:composition-surjective]{\textit{Go to proof.}}
\end{proposition}

\begin{remark*}[Standard quantified statement]
\[
  \operatorname{Surjective}(f)\wedge\operatorname{Surjective}(g)
  \;\Rightarrow\;
  \forall c\in C,\;\exists a\in A: g(f(a))=c.
\]
\end{remark*}

\begin{remark*}[Definition predicate reading]
\[
  \operatorname{Surjective}(g\circ f, C)
  \;\colonequiv\;
  \forall c\in C,\;\exists a\in A: (g\circ f)(a)=c.
\]
\end{remark*}

\NoLocalDependencies
% ---------------------------------------------------------
% Corollary — Composition of Bijections
% ---------------------------------------------------------

\begin{corollary}[Composition of Bijections Is Bijective]
\label{cor:composition-bijective}
Let $f:A\to B$ and let $g:B\to C$. If $f$ is bijective and $g$ is
bijective, then $g\circ f:A\to C$ is bijective.

\smallskip\noindent
\hyperref[prf:composition-bijective]{\textit{Go to proof.}}
\end{corollary}

\begin{remark*}[Standard quantified statement]
\[
\begin{aligned}
&\Bigl[
  \bigl(\forall a_1,a_2\in A\;(f(a_1)=f(a_2)\Rightarrow a_1=a_2)\bigr)
  \wedge
  \bigl(\forall b\in B\;\exists a\in A\;(f(a)=b)\bigr)
 \Bigr] \\
&\quad\wedge
 \Bigl[
  \bigl(\forall b_1,b_2\in B\;(g(b_1)=g(b_2)\Rightarrow b_1=b_2)\bigr)
  \wedge
  \bigl(\forall c\in C\;\exists b\in B\;(g(b)=c)\bigr)
 \Bigr] \\
&\quad\Rightarrow
 \Bigl[
  \bigl(\forall a_1,a_2\in A\;(g(f(a_1))=g(f(a_2))\Rightarrow a_1=a_2)\bigr)
  \wedge
  \bigl(\forall c\in C\;\exists a\in A\;(g(f(a))=c)\bigr)
 \Bigr].
\end{aligned}
\]
\end{remark*}

\begin{remark*}[Definition predicate reading]
\[
  \operatorname{Bijective}(f)\wedge\operatorname{Bijective}(g)
  \;\Longrightarrow\;
  \operatorname{Bijective}(g\circ f).
\]
\end{remark*}

\begin{remark*}[Negated quantified statement]
\[
\begin{aligned}
&\exists a_1,a_2\in A\;\bigl(a_1\neq a_2\;\wedge\;g(f(a_1))=g(f(a_2))\bigr) \\
&\quad\vee\;
  \exists c\in C\;\forall a\in A\;\bigl(g(f(a))\neq c\bigr).
\end{aligned}
\]
\end{remark*}

\begin{remark*}[Negation predicate reading]
\[
  \neg\,\operatorname{Bijective}(g\circ f).
\]
\end{remark*}

\begin{remark*}[Failure modes]
A composite bijection can fail only by losing injectivity or by missing
some target value in $C$.
\end{remark*}

\begin{remark*}[Failure mode decomposition]
\[
\begin{aligned}
\neg\,\operatorname{Bijective}(g\circ f)
&\;\Longleftrightarrow\;
\neg\,\operatorname{Injective}(g\circ f)
\;\vee\;
\neg\,\operatorname{Surjective}(g\circ f).
\end{aligned}
\]
\end{remark*}

\begin{remark*}[Interpretation]
A bijection is an injective and surjective function. Since both of those
properties are preserved by composition, bijectivity is preserved too.
\end{remark*}

\NoLocalDependencies
% ---------------------------------------------------------
% Proposition — Inverse of Bijection
% ---------------------------------------------------------

\begin{proposition}[Inverse of a Bijection Is a Bijection]
\label{prop:inverse-bijection}
$\operatorname{Bijective}(f)\Rightarrow\operatorname{Bijective}(f^{-1})$.

\smallskip\noindent
\hyperref[prf:inverse-bijection]{\textit{Go to proof.}}
\end{proposition}

\begin{remark*}[Standard quantified statement]
\[
  \operatorname{Bijective}(f)
  \;\Rightarrow\;
  \operatorname{InverseFunction}(f, f^{-1}),\;
  f^{-1}\circ f=\mathrm{id}_A,\;
  f\circ f^{-1}=\mathrm{id}_B.
\]
\end{remark*}

\begin{remark*}[Definition predicate reading]
\[
  \operatorname{InverseFunction}(f, f^{-1})
  \;\colonequiv\;
  f^{-1}\circ f = \mathrm{id}_A \;\wedge\; f\circ f^{-1}=\mathrm{id}_B.
\]
\end{remark*}

\NoLocalDependencies
% ---------------------------------------------------------
% Proposition — Preimage Distributes
% ---------------------------------------------------------

\begin{proposition}[Preimage Distributes over Union and Intersection]
\label{prop:preimage-union-intersection}
Let $f:A\to B$, $S,T\subseteq B$. Then:
\begin{enumerate}[label=(\roman*)]
  \item $f^{-1}(S\cup T)=f^{-1}(S)\cup f^{-1}(T)$.
  \item $f^{-1}(S\cap T)=f^{-1}(S)\cap f^{-1}(T)$.
  \item $f^{-1}(S^c)=(f^{-1}(S))^c$.
\end{enumerate}

\smallskip\noindent
\hyperref[prf:preimage-union-intersection]{\textit{Go to proof.}}
\end{proposition}

\begin{remark*}[Standard quantified statement]
\[
\begin{aligned}
(i)\quad &x\in f^{-1}(S\cup T)\iff f(x)\in S\vee f(x)\in T
         \iff x\in f^{-1}(S)\cup f^{-1}(T). \\
(ii)\quad &x\in f^{-1}(S\cap T)\iff f(x)\in S\wedge f(x)\in T
          \iff x\in f^{-1}(S)\cap f^{-1}(T).
\end{aligned}
\]
\end{remark*}

\begin{remark*}[Negated quantified statement]
\[
  x\notin f^{-1}(S)
  \;\iff\; f(x)\notin S
  \;\iff\; x\in f^{-1}(S^c).
\]
\end{remark*}

\begin{remark*}[Interpretation]
Preimage is a complete lattice homomorphism: it preserves all set
operations including complement. Contrast with image: $f(A\cap B)
\subseteq f(A)\cap f(B)$ with strict containment possible when $f$
is not injective. Preimage identities are load-bearing in topology
(preimages of open sets are open) and measure theory (preimages of
measurable sets are measurable).
\end{remark*}

\NoLocalDependencies
